\section{Video Preprocessing Algorithms}
\label{sec:preprocessing}

\subsection{ASR Transcription (Task 1)}

The first preprocessing task extracts audio from the uploaded video and transcribes it using Alibaba DashScope's Qwen3-ASR service.

\subsubsection{Audio Extraction}

Audio is extracted using FFmpeg with the following parameters:

\begin{lstlisting}[language=bash, caption={Audio Extraction Command}]
ffmpeg -i input.mp4 -vn -acodec pcm_s16le \
       -ar 16000 -ac 1 output.wav
\end{lstlisting}

The audio is converted to 16kHz mono PCM WAV format, which is the optimal format for the ASR service. The extracted audio file is uploaded to Tencent Cloud Object Storage (COS) for temporary hosting during transcription.

\subsubsection{Qwen3-ASR Transcription}

We use the asynchronous file transcription API of Qwen3-ASR, which provides:

\begin{itemize}
    \item \textbf{Sentence-level timestamps:} Each transcribed sentence includes precise begin and end times
    \item \textbf{Language detection:} Automatic identification of spoken language (useful for multilingual lectures)
    \item \textbf{Emotion tagging:} Detection of speaker emotional states (neutral, happy, excited, etc.)
    \item \textbf{Punctuation restoration:} Proper punctuation insertion for readable transcripts
\end{itemize}

The ASR service processes the audio file and returns results in JSON format:

\begin{lstlisting}[language=json, caption={ASR Response Format}]
{
  "sentences": [
    {
      "begin_time": 0.0,
      "end_time": 4.5,
      "text": "Welcome to today's lecture on machine learning.",
      "language": "en",
      "emotion": "neutral"
    },
    ...
  ]
}
\end{lstlisting}

Transcription typically completes in 1-2 minutes for a 60-minute lecture, depending on audio quality and speaking rate.

\subsection{HLS Encoding (Task 2)}

HLS (HTTP Live Streaming) encoding prepares videos for adaptive bitrate streaming in the frontend player.

\subsubsection{Multi-Resolution Transcoding}

The video is transcoded into four resolution variants:

\begin{table}[H]
    \centering
    \caption{HLS Encoding Resolutions}
    \begin{tabular}{@{}lll@{}}
        \toprule
        \textbf{Name} & \textbf{Resolution} & \textbf{Bitrate} \\ \midrule
        1080p & 1920$\times$1080 & 5000 kbps \\
        720p & 1280$\times$720 & 2800 kbps \\
        480p & 854$\times$480 & 1400 kbps \\
        360p & 640$\times$360 & 800 kbps \\ \bottomrule
    \end{tabular}
\end{table}

Each variant is segmented into 10-second TS (MPEG Transport Stream) chunks. A master M3U8 playlist references all variant playlists, enabling the frontend player to automatically select the optimal quality based on available bandwidth.

\subsubsection{FFmpeg Command Structure}

\begin{lstlisting}[language=bash, caption={HLS Encoding FFmpeg Command}]
ffmpeg -i input.mp4 \
  -map 0:v:0 -map 0:a:0 \
  -s 1920x1080 -b:v 5000k -c:v libx264 \
  -hls_time 10 -hls_playlist_type vod \
  -hls_segment_filename "streams/video_1080p_%03d.ts" \
  "streams/video_1080p.m3u8"
# (Repeated for each resolution)
\end{lstlisting}

\subsection{SSIM Slide Detection (Task 3)}

The SSIM slide detection algorithm identifies slide transitions by analyzing frame-to-frame visual changes. This is a core innovation of our system, enabling accurate segmentation without manual annotation.

\subsubsection{Structural Similarity Index}

The Structural Similarity Index (SSIM) \cite{wang2004image} measures the perceived visual similarity between two images. Unlike simple pixel-wise difference metrics, SSIM accounts for luminance, contrast, and structural information, making it more aligned with human visual perception.

For two image patches $x$ and $y$, SSIM is computed as:

\begin{equation}
    \text{SSIM}(x, y) = \frac{(2\mu_x\mu_y + c_1)(2\sigma_{xy} + c_2)}{(\mu_x^2 + \mu_y^2 + c_1)(\sigma_x^2 + \sigma_y^2 + c_2)}
\end{equation}

where:
\begin{itemize}
    \item $\mu_x, \mu_y$ are the mean intensities
    \item $\sigma_x^2, \sigma_y^2$ are the variances
    \item $\sigma_{xy}$ is the covariance
    \item $c_1, c_2$ are stabilization constants
\end{itemize}

SSIM values range from 0 (completely different) to 1 (identical). A drop below the threshold (default 0.7) indicates a slide change.

\subsubsection{Multithreaded Processing Pipeline}

Listing~\ref{alg:ssim-detection} details our multithreaded SSIM detection approach.

\begin{lstlisting}[language=Python, caption={Multithreaded SSIM Slide Detection}, label={alg:ssim-detection}]
Input: Video file V, threshold tau, sampling rate fs
Output: List of slide change timestamps T

T = empty set
last_change = 0
open video V with OpenCV
prev_frame = read and preprocess first frame
create thread pool with N=16 workers

while frames remaining:
    curr_frame = read and preprocess next frame
    if frame interval >= 1/fs:
        submit SSIM computation to thread pool
        ssim = SSIM(prev_frame, curr_frame)
        if ssim < tau and t - last_change >= delta_t_min:
            T = T union {t}
            last_change = t
        prev_frame = curr_frame

return sorted(T)
\end{lstlisting}

\subsubsection{Preprocessing Steps}

Each frame undergoes preprocessing before SSIM computation:

\begin{enumerate}
    \item \textbf{Resize:} Scale to 240 pixels width (height proportional) to reduce computational load
    \item \textbf{Grayscale Conversion:} Convert from RGB to single-channel grayscale
    \item \textbf{Normalization:} Scale pixel values to [0, 1] range
\end{enumerate}

These preprocessing steps reduce each frame from approximately 2 megapixels (1920$\times$1080 RGB) to about 32,000 pixels (240$\times$135 grayscale), achieving a 60$\times$ reduction in data size while preserving structural information necessary for slide change detection.

\subsubsection{Performance Optimization}

Several optimizations enable efficient processing:

\begin{itemize}
    \item \textbf{Frame Sampling:} Processing at 10 FPS instead of native 30 FPS reduces frame count by 3$\times$ while maintaining detection accuracy (slide changes are typically sustained for multiple frames)
    \item \textbf{Thread Pool:} 16 worker threads parallelize SSIM computation while video decoding remains single-threaded (OpenCV requirement)
    \item \textbf{Cooldown Period:} Minimum 5-second interval between detections prevents false positives from animations or video transitions
    \item \textbf{Memory Efficiency:} Only two frames are held in memory per worker, resulting in ~1MB peak memory usage for the thread pool
\end{itemize}

\begin{table}[H]
    \centering
    \caption{SSIM Detection Performance}
    \begin{tabular}{@{}lll@{}}
        \toprule
        \textbf{Video Length} & \textbf{Frames @ 10 FPS} & \textbf{Processing Time} \\ \midrule
        10 min & ~6,000 & 10-20 s \\
        30 min & ~18,000 & 30-50 s \\
        60 min & ~36,000 & 60-120 s \\
        120 min & ~72,000 & 2-4 min \\ \bottomrule
    \end{tabular}
\end{table}

\subsection{Thumbnail Generation (Task 4)}

Once slide change timestamps are identified, representative thumbnails are extracted at each transition point.

\subsubsection{Extraction Algorithm}

For each detected slide change at timestamp $t$:

\begin{enumerate}
    \item Seek to frame at time $t + 0.5$ seconds (slight offset to ensure new slide is fully visible)
    \item Extract frame and resize to 640$\times$360 pixels
    \item Save as JPEG with 85\% quality
    \item Store metadata (timestamp, file path) in database
\end{enumerate}

These thumbnails serve multiple purposes: visual preview in the transcript viewer, representative images for sections, and visual context for LLM-based knowledge extraction.
